\chapter{State of the Art and Technical Foundations}

\section{Introduction}

JANUS combines technologies that are often deployed independently. Honeytokens
provide low-noise tripwires, host auditing records local activity, a SIEM/HIDS
normalizes evidence, and generative AI varies the decoy content. This chapter
explains the role and limits of each component.

\section{Cyber deception and honey documents}

Cyber deception deliberately exposes false information or services in order to
influence an adversary and improve defender observation. In contrast with a
production resource, a properly isolated decoy should have no routine use. This
raises the investigative value of an interaction, but it does not automatically
make every event malicious: antivirus, indexing software, backup agents, and
authorized tests can also touch a file. A mature design therefore records
context and supports suppression rather than assuming intent from one event.

A honey document is a decoy file with a plausible business purpose and one or
more observation mechanisms. Its realism depends on its filename, internal
structure, metadata, placement, and consistency with neighboring files. Its
safety depends on the absence of real secrets and dangerous active content.

\section{Honeytokens and remote callbacks}

Canarytokens implement lightweight tripwires that notify a defender when a
unique token is used. The official getting-started guide describes the basic
model as create, place, and wait for an alert \cite{canaryGettingStarted}.
Activation is token-specific.

\subsection{Office tokens}

Microsoft Word and Excel tokens can alert when opened in their native
applications \cite{canaryOffice}. Non-macro variants rely on a network lookup
for embedded external content. This is preferable for JANUS because it avoids
asking users to enable macros. It is nevertheless conditional: protected view,
external-content policy, offline operation, or network filtering can prevent
the request.

\subsection{Web breadcrumbs}

A URL embedded in CSV, JSON, YAML, ENV, or ZIP content is passive. Opening the
file in a text editor does not execute a network request. The alert occurs only
if a person or tool uses the URL. JANUS therefore labels this mode as
user-action-required instead of automatic.

\subsection{AWS key tokens}

An AWS API Keys Canarytoken is a decoy access-key pair. Thinkst documents that
an alert is generated when the keys are used through the AWS API and notes a
possible delay of 2 to 30 minutes through AWS logging infrastructure
\cite{canaryAws}. Merely displaying an ENV file must not be reported as a
remote token activation.

\section{Windows SACL and event 4663}

A System Access Control List specifies which access attempts on a securable
object must be audited. It is distinct from a discretionary ACL: the SACL does
not grant access; it requests audit records when selected rights are exercised.

Windows Security event 4663 states that an operation was performed on an
object. Microsoft explains that it is generated only when the object's SACL
contains the required audit entry and that, unlike event 4656, it represents a
right that was used rather than merely requested \cite{microsoft4663}. Useful
fields include:

\begin{itemize}
  \item \texttt{ObjectName}: the accessed file or directory;
  \item \texttt{ProcessName} and \texttt{ProcessId}: the local process;
  \item \texttt{SubjectUserName}, domain, SID, and logon ID: the security
  context;
  \item \texttt{AccessMask} and \texttt{AccessList}: the exercised rights;
  \item timestamp, computer name, and event record identifier.
\end{itemize}

Event 4663 is strong evidence that a right was exercised on the monitored host,
but it is not a semantic record of user intention. A read can result from Word,
PowerShell, an antivirus engine, an indexer, or a copy utility. A remote
download requires corroborating SMB, authentication, process, or network
events.

\section{Wazuh file and event monitoring}

Wazuh provides agent-based collection, rules, indexing, and visualization.
Its File Integrity Monitoring module tracks file creation, modification, and
deletion by comparing checksums and attributes \cite{wazuhFim}. In who-data
mode on Windows, Wazuh uses the Microsoft audit subsystem and can configure the
required audit policy and SACL \cite{wazuhWhodata}. JANUS additionally consumes
the Security event channel and applies dedicated rules to audited access under
the project deployment root.

For Linux, the kernel audit system can monitor read, write, execute, and
attribute access using keyed audit rules; Wazuh provides a documented path for
collecting and visualizing these events \cite{wazuhAuditd}. JANUS includes an
adapter for Linux audit signals, but the validated demonstration environment is
Windows.

\section{Generative AI for deception content}

An LLM can produce varied business narratives and reduce the repetitive nature
of static templates. It should not directly control deployment, token creation,
or filesystem paths. JANUS therefore uses the model for constrained text only.
Deterministic code selects the format, creates the token, assembles the file,
validates its structure, chooses the destination, and records evidence. A local
fallback keeps the project testable when a remote model is not configured.

The main risks are leakage of secrets in prompts, uncontrolled model output,
hallucinated operational claims, prompt injection embedded in source material,
and non-reproducible content. The implementation addresses these risks through
prompt minimization, a controlled corpus, coherence checks, forbidden-pattern
scanning, and post-generation artifact validation.

\section{Comparative analysis}

\begin{table}[H]
\centering
\caption{Comparison of observation mechanisms used by JANUS}
\label{tab:mechanism-comparison}
\small
\begin{tabular}{p{0.19\linewidth}p{0.24\linewidth}p{0.24\linewidth}p{0.24\linewidth}}
\toprule
Mechanism & Direct evidence & Main value & Main limitation \\
\midrule
Windows 4663 & Account, process, object, access right, time & Local interaction context & Only on an audited host; read semantics are limited \\
Wazuh & Normalized alert, rule, agent, indexed evidence & Collection and correlation & Depends on agent, rules, and indexer availability \\
Office token & Remote HTTP/DNS activation and request metadata & Detection after the file leaves the host & Can be blocked by Office or network policy \\
AWS keys token & Use of a decoy key against AWS APIs & High-value credential-use signal & Opening the file does not activate it; alert may be delayed \\
Web breadcrumb & Request to a unique URL & Simple cross-platform tripwire & Requires a deliberate click or tool action \\
\bottomrule
\end{tabular}
\end{table}

\section{Positioning of JANUS}

JANUS is not a replacement for Wazuh or Canarytokens. Its contribution is the
orchestration layer between them: it generates and registers each decoy,
selects the appropriate token, controls deployment, correlates endpoint events
with the active registry, and explains sensor limits in one interface. The
design is intentionally low interaction and defensive.

\section{Conclusion}

No single sensor provides every desired attacker attribute. The architecture
must combine endpoint evidence and remote tripwires while preserving the
different confidence levels of each source. Chapter 3 turns these principles
into system requirements and components.

